\section{Introduction}
\label{sec:introduction}

The electroweak sector contains a dimensionless pole-spectrum ratio,
\begin{equation}
  \sPole = 1-\frac{M_{W,\rm pole}^2}{M_{Z,\rm pole}^2},
  \label{eq:pole-angle-intro}
\end{equation}
which is fixed experimentally by the pole positions of the charged and neutral weak vector bosons.  In the Standard Model at tree level the same ratio is written as an angular datum in coupling space, while the Higgs vacuum expectation value fixes the common radial scale of the vector masses.  This separation between projective and radial data gives the working language of the present paper.

The Rivero--de Vries observation starts from a mass operator built from Poincare Casimir data and produces a dimensionless two-branch equation.  In the normalized notation used below, the equation is
\begin{equation}
  x^2+Jx-J=0,
  \qquad J=s(s+1).
  \label{eq:intro-secular}
\end{equation}
The positive branch gives
\begin{equation}
  \sdV = 1-\frac{x_+(3/4)}{x_+(2)}=0.22310132\ldots .
  \label{eq:intro-devries}
\end{equation}
The original note emphasized the proximity of this number to the mass-shell weak angle and motivated the operator through Casimir invariants and high-spin behavior \cite{Rivero2006}.  The present manuscript treats the numerical proximity as descriptive.  It records a structural clue.  Eq.~\eqref{eq:intro-devries} has clue status in this draft.  The central question asks which physical object carries the square root in Eq.~\eqref{eq:intro-secular}, and why that object is read at the low-energy W/Z pole ratio.

This placement question is the conceptual core.  A standard grand-unified lineage uses representation theory to obtain a clean weak-angle value at a high scale.  In SU(5), and in Spin(10) with the usual fermion content, the famous group-theoretic value is \(\sin^2\theta_W=3/8\); contact with the observed low-energy angle is then made through running and thresholds \cite{Britto2021}.  The Rivero--de Vries construction reverses the burden.  Its clean value is close to the pole-mass ratio itself.  Therefore the manuscript must explain why a spectral datum should attach to the pole spectrum.  A unification-scale boundary condition would define a different reading of the same algebra.

The string-theory posture uses the older dual-model and Regge ambition as part of the physical context.  Open strings, endpoint data, Kaluza--Klein spectra, and boundary conditions are treated as possible mechanisms for organizing low-energy particle structure.  String theory also incorporates Kaluza--Klein ideas, grand-unified representations, and supersymmetric structure in transmuted form \cite{Polchinski1994}.  The working problem is to ask whether those tools can produce an electroweak secular datum at the observed pole spectrum.  The numerical clue remains a clue until that mechanism is written.

The construction has four layers.  The first layer is algebraic: a two-channel secular equation has roots whose ratio is fixed by the Casimir variable.  The second layer is electroweak: the W/Z mass ratio is a projective datum along the broken-to-unbroken ray of the gauge-Higgs system.  The third layer is spectral: the square root must be assigned to a pole, running, compactification, or boundary-condition object.  The fourth layer is dynamical: strings, branes, Kaluza--Klein boundary conditions, or singular \(G_2\) geometries may generate the secular equation or its branch structure.

The status of the argument is deliberately stratified.  The algebraic calculation is explicit.  The pole-mass comparison is a descriptive clue whose source audit remains a separate task.  The assignment of \(J=3/4\) to the charged W comparison and \(J=2\) to the neutral Z comparison is the central open analytical calculation unless it is derived in Sec.~\ref{sec:ew-embedding}.  The negative branch is retained because it may encode the scalar/order-parameter side of electroweak symmetry breaking.

The paper is organized as follows.  Section~\ref{sec:casimir} gives the secular operator and its roots.  Section~\ref{sec:pole} defines the pole observable and the placement problem.  Section~\ref{sec:ew-embedding} formulates the projective electroweak interpretation.  Section~\ref{sec:negative} records the negative branch and its possible Higgs-scale role.  Sections~\ref{sec:string-kk-g2}--\ref{sec:route-comparison} develop the string, endpoint, Kaluza--Klein, and \(G_2\) mechanism tests in a common two-channel language.  Sections~\ref{sec:flavour} and \ref{sec:global-form} delimit the flavor and global-gauge-group issues.  Section~\ref{sec:status-tests} gives the status of the analytical program and its conceptual tests.

\paragraph*{Section status.}  This introduction states the conceptual placement problem and labels the central unproved step.
